\documentclass[answers,12pt,addpoints]{exam} 
\usepackage{import}

\import{C:/Users/prana/OneDrive/Desktop/MathNotes}{style.tex}

% Header
\newcommand{\name}{Pranav Tikkawar}
\newcommand{\course}{Stats 596}
\newcommand{\assignment}{Notes}
\author{\name}
\title{\course \ - \assignment}

\begin{document}
\maketitle
\tableofcontents
\newpage

\section*{Preface \& Comments }

\section{Chapter 1}
\subsection{Lecture 1: 9/1}
\begin{remark}
    [Data of Model]
    $(X_i, Y_i), i=1,2,\ldots,n$ \quad $X_i$ is the covariate, $Y_i$ is the response. \\

    $$ Y_i = X_i^T \beta^* + \epsilon_i, i = 1,2,\ldots,n$$
    where $\beta^* \in \mathbb{R}^p$ is an unknown coeeficent vecotr. and $\epsilon_i$ is the noise term. 
\end{remark}

\begin{definition}
    [Fixed Design.]
    $X_1, X_2, \ldots, X_n$ are fixed and known and $Y_1, Y_2, \ldots, Y_n$ are random. More speifically $Y_i$ are independent and $E[Y_i] = X_i^T \beta^*$ 

    This is the classical linear regression model. And ignore the randomness in $X_i$'s by conditioning some on $X_i$'s.

    Most of the theory of historical linear regression is based on the fixed design assumption.
\end{definition}
\begin{definition}
    [Random Design]
    $(X_i, Y_i), i=1,2,\ldots,n$ are independent and identically distributed (i.i.d.) random variables, note they are pairs of random variables.
    $$Y_i|X_i \sim N(X_i^T \beta^*, \sigma^2)$$
    $$\epsilon_i|X_i \sim N(0, \sigma^2)$$
    where $\epsilon_i = Y_i - X_i^T \beta^*$ is the noise term.
    More specifically $(X_i, Y_i) \sim P_{XY}$, where $P_{XY}$ is the joint distribution of $(X,Y)$ and $E[Y|X] = X^T \beta^*$

    $E[Y|X] = X^T \beta^*$ is the conditional expectation of $Y$ given $X$ and is a function of $X$.
    $E[\epsilon|X] = 0$ is the conditional expectation of $\epsilon$ given $X$ and is a function of $X$.
\end{definition}
\begin{definition}
    [Homoscedasticity noise]
    Fixed: $\Var(y_i) = \Var(\epsilon_i) = \sigma^2$
    Random: $\Var(Y_i|X_i) = \sigma^2$ 
\end{definition}

\subsection{Lecture 2: 9/3}
\begin{remark}
    [Estimation: LS]
    $$\hat{\beta} = \arg \min_{\beta} \sum_{i=1}^{n} (Y_i - X_i^T \beta)^2$$
    This is the least squares estimator for the linear regression model.

    \textbf{Closed-Form Solution}
    The closed-form solution for the least squares estimator is:
    $$\hat{\beta} = (X^T X)^{-1} X^T Y$$
    where $X$ is the design matrix and $Y$ is the response vector.
\end{remark}
\begin{remark}
    [Estimatior of $\sigma^2$]
    Under Homoscedasticity, the estimator of $\sigma^2$ is:
    $$\hat{\sigma}^2 = \frac{1}{n-p} \sum_{i=1}^{n} (Y_i - X_i^T \hat{\beta})^2$$
    where $p$ is the number of parameters in the model.

    Under MLE, the estimator of $\sigma^2$ is:
    $$\hat{\sigma}_\text{MLE}^2 = \frac{1}{n} \sum_{i=1}^{n} (Y_i - X_i^T \hat{\beta})^2$$
    where $p$ is the number of parameters in the model. Remember that this requires a normality assumption on the noise term $\epsilon_i$.
\end{remark}
\begin{proposition}
    [Homoscedastic noise in fixed design]
    \begin{align*}
        E[\hat{\beta}] = \beta^*\\
        \var(\hat{\beta}) = (n\hat{H})^{-1}{\sigma^*}^2\\
        E[\hat{\sigma}^2] = {\sigma^*}^2
    \end{align*}
    Where $\hat{H} = \frac{1}{n} \sum_{i=1}^{n} X_i X_i^T$ is the empirical covariance matrix of the covariates $X_i$'s.
\end{proposition}
\begin{remark}
    $E[\hat{\beta}] = \beta^*$
    \begin{proof}
        \begin{align*}
            E[\hat{\beta}] &= E[(X^T X)^{-1} X^T Y]\\
            &= (X^T X)^{-1} X^T E[Y]\\
            &= (X^T X)^{-1} X^T E[X \beta^* + \epsilon]\\
            &= (X^T X)^{-1} X^T (X \beta^* + E[\epsilon])\\
            &= (X^T X)^{-1} X^T (X \beta^* + 0)\\
            &= (X^T X)^{-1} X^T X \beta^*\\
            &= \beta^*
        \end{align*}
    \end{proof}
    Note we only use that $E[\epsilon] = 0$ and not the distribution of $\epsilon$
\end{remark}
\begin{definition}
    [Wald Confidence Interval]
    $$\hat{\beta}_j \pm z_{1-\alpha/2} \sqrt{\var(\hat{\beta}_j)}$$
    where $z_{1-\alpha/2}$ is the $(1-\alpha/2)$ quantile of the standard normal distribution.
    This is a confidence interval for the $j$-th component of the parameter vector $\beta$.
    For $\hat{\mu}$ the wald CI is $\hat{\mu} \pm z_{1-\alpha/2} \sqrt{\var(\hat{\mu})}$ where $\hat{\mu}$ is the estimator of the mean $\mu$.
    
    Point estimator and Variance Estimator (asymptotic normality) there are hings to remember for the wald CI 

\end{definition}

\subsection{Lecture 3: 9/10}
\begin{remark}
    Remember Lin Reg:
    Fixed Design: $Y_i = X_i^T \beta^* + \epsilon_i, i = 1,2,\ldots,n$ where $\epsilon_i$ are independent and $E[\epsilon_i] = 0$ and $\var(\epsilon_i) = \sigma^2$. \\
    CI for $\beta_j^*$ with homoscedasticity: $$\hat{\beta}_j \pm z_{1-\alpha/2} \hat\sigma \sqrt{\frac{\hat{H}^{-1}_{jj}}{n}}$$
\end{remark}

\begin{proposition}
    [Heteroscedastic noise]
    \begin{align*}
        E[\hat{\beta}] &= \beta^*\\
        \var(\hat{\beta}) &= \hat{H}^{-1} \left(\frac{1}{n^2} \sum_{i} {\sigma_i^{*}}^2 X_i X_i^T \right) \hat{H}^{-1}\\
        \text{where} \quad {\sigma_i^{*}}^2 &= \var(\epsilon_i) = \var(Y_i|X_i)\\
    \end{align*}
    if $\sigma_i^{*}$ are constant then we can pull it out and it is clear that it matches the homoscedastic case.
    $$ \hat{\Sigma} = \hat{H}^{-1} \frac{1}{n} \left(\sum_{i=1}^{n} (Y_i - X_i^T \hat{\beta})^2 X_i X_i^T\right) \hat{H}^{-1}$$
\end{proposition}

\begin{remark}
    \begin{align*}
        \hat{\Sigma}(\beta^*) &= \hat{H}^{-1} \frac{1}{n} \left(\sum_{i=1}^{n} (Y_i - X_i^T \beta^*)^2 X_i X_i^T\right) \hat{H}^{-1}\\
        \sqrt{n}(\hat{\beta} - \beta^*) &\to_ D N(0, \hat{\Sigma}(\beta^*))\\
        ||\hat{\Sigma}(\hat{\beta}) - \hat{\Sigma}(\beta^*)||_F &\to 0 \quad \text{in probability}
    \end{align*}
\end{remark}

\begin{remark}
    [Random Design]
    $\hat{\beta}$ LS estimator is stil good

    "Conditional inference": Given $X_1, X_2, \ldots, X_n$, then fixed design\\
    "Unconditional inference": $X_1, X_2, \ldots, X_n$ are random, then random design: 
    $$E(\hat{\beta}) = E[E[\hat{\beta}|X_i]] = E[\beta^*] = \beta^*$$
    $$\var(\hat{\beta}) = E[\var(\hat{\beta}|X_i)] + \var(E[\hat{\beta}|X_i]) = E[\var(\hat{\beta}|X_i)]$$
\end{remark}
\begin{proposition}
    \begin{align*}
        \sqrt{n}(\hat{\beta} - \beta^*) &\to_D N(0, \Sigma(\beta^*))\\
        \hat{\Sigma}(\hat{\beta}^*) &\to_P \Sigma(\beta^*)\\
        \text{where} \quad \Sigma(\beta^*) &= H^{-1}G(\beta^*) H^{-1}\\
        \text{and} \quad H &= E[X_i X_i^T]\\
        \text{and} \quad G(\beta^*) = E[(Y_i - X_i^T \beta^*)^2 X_i X_i^T]\\
    \end{align*}
    \begin{proof}
        \begin{align*}
            \hat \beta = \tilde{E}^{-1}[XX^T] \tilde{E}[XY]\\
            \hat \beta - \beta^* = \tilde{E}^{-1}[XX^T] \tilde{E}[XY] - \tilde{E}^{-1}[XX^T] \tilde{E}[XX^T] \beta^*\\
            = \tilde{E}^{-1}[XX^T] \tilde{E}[X(Y - X^T \beta^*)]\\
            \sqrt{n}(\hat \beta - \beta^*) = \tilde{E}^{-1}[XX^T] \sqrt{n} \tilde{E}[X(Y - X^T \beta^*)]\\
            \text{By CLT, } \sqrt{n} \tilde{E}[X(Y - X^T \beta^*)] \to_D N(0, G(\beta^*))
        \end{align*}
    \end{proof}
    Where $\tilde{E}$ is the empirical expectation operator.
\end{proposition}

\begin{theorem}
    [Slutsky's Theorem]
    If $X_n \to_D X$ and $Y_n \to_P c$ then $X_n + Y_n \to_D X + c$ and $X_n Y_n \to_D cX$ and $X_n / Y_n \to_D X / c$ if $c \neq 0$
\end{theorem}

\begin{remark}
    [Application of Slutsky's Theorem]
    We can use this to show $\var(H^{-1} Z) = H^{-1} \var(Z) H^{-1}$ and $\hat{\Sigma}(\hat{\beta}) \to_P \Sigma(\beta^*)$.
\end{remark}

\begin{definition}
    [Partial Linear Regression]
    $$E[Y|X] = X^T \theta^* + g^*(Z)$$
    where $g(Z)$ is an unknown function of $Z$ and $Z$ is another covariate. This is a semi-parametric model.
\end{definition}

\begin{proposition}
    [1.9]
    $\hat \theta$ is doubly frobust for $\theta^*$ if either $g^*(Z)$ is correctly specified or $E[X|Z]$ is correctly specified.
\end{proposition}

\begin{proposition}
    Let $V(\theta^*) = \frac{E\{\phi^2(z,w,y;\theta*, \bar \alpha, \bar \gamma) \}}{E^2\{Z (Z-W^T\bar \gamma)\}}$
    If either g or f model is correct then $\sqrt{n}(\hat \theta - \theta^*) \to_D N(0, V(\theta^*))$ and $\hat V(\hat \theta) \to_P V(\theta^*)$.
\end{proposition}

\begin{remark}
    

    $\frac{1}{n}\sum \{ h(X_i ; \hat \eta) - h(X_i ; \bar \eta)\} \to_P 0$

    This says that, as $n \to \infty$, the sample-average difference between
    evaluating $h$ at the estimated nuisance parameter $\hat\eta$ and at its
    target or probability limit $\bar\eta$ converges to zero in probability.
    Thus, replacing $\bar\eta$ with $\hat\eta$ has a negligible effect on the
    average for large samples. Equivalently, for every $\varepsilon>0$,
    $$
    P\left(\left|\frac{1}{n}\sum_{i=1}^n
    \{h(X_i;\hat\eta)-h(X_i;\bar\eta)\}\right|>\varepsilon\right)\to 0.
    $$
    
    $n^{1/2} (\hat \eta - \bar \eta) = O_P(1)$
    
\end{remark}

\begin{definition}
    $\phi = (Z - W^T \bar \gamma) (Y - W^T \bar \alpha - Z\theta)$

    But if we choose 
    $$\phi_{new} (Z - \frac{\exp(W^T \gamma)}{1+\exp(W^T \gamma)})(Y - W^T \bar \alpha - Z\theta)$$
    then this is still mean zero and is doubly robust.
\end{definition}




\end{document}